% titel: Blatt 2, A3: Binomischer Lehrsatz

% Teil I: Werkzeuge, kein Lösungsweg.

Auswerten einer Binomialsumme, Wahrscheinlichkeiten bei DNA-Replikation
und Nachweis, dass die Mutationswahrscheinlichkeiten zusammen $1$ ergeben.

\begin{werkzeugkasten}[Binomischer Lehrsatz]
\begin{equation}
  (a+b)^n=\sum_{k=0}^{n}\binom{n}{k}a^k b^{n-k}.
\end{equation}
Setzt man $a=-2$, $b=1$: $\sum_{k=0}^{n}\binom{n}{k}(-2)^k$
$=(-2+1)^n=(-1)^n$, also $1$ bzw. $-1$ je nach Parität von $n$.
Setzt man $a=p$, $b=1-p$: $\sum_{k=0}^{n}\binom{n}{k}p^k (1-p)^{n-k}=1$.
\end{werkzeugkasten}

\begin{werkzeugkasten}[Binomialverteilung]
Bei $n$ unabhängigen Versuchen mit Erfolgswahrscheinlichkeit $p$:
\begin{equation}
  \mathrm P(\text{genau }k\text{ Erfolge})=\binom{n}{k}p^k(1-p)^{n-k}.
\end{equation}
``Keine Mutation'' entspricht $k=0$: $\,(1-p)^{n}$. Der Lehrsatz mit
$a=p,b=1-p$ zeigt, dass alle $\binom{n}{k}p^k(1-p)^{n-k}$ zusammen $1$
ergeben -- sie bilden eine echte Wahrscheinlichkeitsverteilung.
\end{werkzeugkasten}

\begin{achtung}[Darauf achten]
\begin{itemize}
  \item Den Vorzeichen-Fall: $(-2)^k$ wechselt das Vorzeichen mit der
        Parität; das Gesamtergebnis $(-1)^n$ hängt nur von $n\bmod2$ ab.
  \item ``Keine Mutation'': genau $k=0$ einsetzen, nicht $1-p^n$.
  \item Beim Lehrsatz-Beweis die richtigen $a,b$ wählen: für
        $\sum p^k(1-p)^{n-k}$ braucht man $a=p$, $b=1-p$, da $p+(1-p)=1$.
\end{itemize}
\end{achtung}
